\addtocontents{toc}{\protect\newpage}
\chapter{Zusammenfassung und Schlussfolgerungen}
\label{ch:zusammenfassung}

Die Arbeit hat untersucht, wie Prozessineffizienzen in \gls{erp}-Systemen mit Process Mining evidenzbasiert aufgedeckt werden können, anhand eines synthetischen, vollständig reproduzierbaren \gls{sap}-\gls{o2c}-Event-Logs mit bekannter Grundwahrheit. Zu F1 ist die zentrale Datenstruktur das Event-Log aus Fällen, Aktivitäten und Zeitstempeln, standardisiert im \gls{xes}-Format. In \gls{sap} \gls{s/4hana} liegt es nicht nativ vor, lässt sich aber aus den Belegtabellen (\gls{vbak}/\gls{vbap}, \gls{likp}/\gls{lips}, \gls{vbrk}/\gls{vbrp}, \gls{bkpf}/\gls{bseg}), der Belegflusstabelle \gls{vbfa} und den Änderungsbelegen \gls{cdhdr}/\gls{cdpos} rekonstruieren. Kritisch sind die Wahl der Case-ID, die Zeitstempelgranularität und die Datenqualität. Zu F2 scheiterte der Alpha-Algorithmus im empirischen Vergleich auf 12.000 Fällen an Schleifen und Rauschen (Fitness 0,646), während Heuristics Miner (0,972) und insbesondere Inductive Miner (0,996 bei Precision 0,913) brauchbare Modelle lieferten. Für den Soll-Ist-Abgleich erwiesen sich Alignments als notwendig, weil einfaches Token-Replay Abweichungsaktivitäten außerhalb des Soll-Modells übersieht. Zu F3 hat die Analyse alle sechs injizierten Störungsmuster ohne Vorwissen wiedergefunden und quantifiziert, von den Kreditsperren (34,1 \% der Fälle) bis zu den Gutschriften mit Nachlieferung (4,4 \%, Tabelle~\ref{tab:ergebnis-fitgap}). Nur 38,9 Prozent der Fälle liefen konform zum \gls{sap}-Referenzprozess.

Die abgeleiteten Maßnahmen bleiben vollständig im \gls{sap}-Standard und adressieren jede gemessene Abweichung mit einem konkreten Customizing- oder Organisationshebel. Eine Ex-post-Bewertung ihrer Effektivität entfällt studiendesignbedingt, weil die Maßnahmen abgeleitet und nicht umgesetzt wurden. An ihre Stelle tritt das Kennzahlenset in Anhang A3 mit Ist- und Zielwerten als Bewertungsmaßstab des Prozessmonitorings. Der Validierungsansatz aus bekannter Grundwahrheit und gemessener Wiederfindung belegt die Tauglichkeit der Methodenkette aus Extraktion, Discovery und Conformance Checking. Einschränkend zeigt der Log modelliertes Verhalten. Reale \gls{erp}-Daten sind unordentlicher und werfen Datenschutzfragen auf, die hier bewusst ausgeklammert wurden. Zudem hängt die Bedeutung einzelner Störungsmuster vom Geschäftskontext ab. In der Fallstudie von Adolfsson/Saleh etwa spielten Kreditsperren im Kaffeegeschäft mit stabiler Nachfrage praktisch keine Rolle \parencite[S.~51]{adolfsson2026o2c}. Die Übertragung auf ein reales System erfordert daher vor allem die in Kapitel~\ref{sec:stamm-bewegungsdaten-sap} beschriebene Extraktionsarbeit, die Algorithmen selbst sind dafür bereit.

Über die drei Forschungsfragen hinaus zeigt die Arbeit einen methodischen Mehrwert gegenüber der klassischen, workshopbasierten Fit/Gap-Analyse. Wo Letztere auf die notorisch unvollständige Dokumentation und die Erinnerung der Beteiligten angewiesen ist, stützt sich Process Mining auf die vollständige, objektive Aufzeichnung des \gls{erp}-Systems. Die Lücken werden nicht mehr geschätzt, sondern gezählt, und die Priorisierung der Maßnahmen erfolgt nach gemessener Häufigkeit und Liegezeit statt nach gefühlter Dringlichkeit. Damit wird die Fit/Gap-Analyse von einer einmaligen, subjektiven Momentaufnahme zu einem wiederholbaren, evidenzbasierten Verfahren, was gerade für das kontinuierliche Monitoring nach dem Go-Live von Bedeutung ist.

Für die Zukunft zeichnen sich drei Entwicklungen ab. Objektzentriertes Process Mining löst die Beschränkung auf eine einzige Case-ID auf und adressiert damit die in Kapitel~\ref{sec:datenstrukturen} diskutierten Konvergenz- und Divergenzprobleme \parencite[S.~3 ff.]{vanderaalst2019objectcentric}. Predictive und Prescriptive Analytics erweitern die Diagnose um Vorhersagen und Handlungsempfehlungen \parencite[Folie~470 f.]{coletta2026_1407}. Dass dieser Schritt marktreif ist, zeigt das \gls{sap}-eigene Auftragsabwicklungsmodell mit Prognosedaten (Process Forecast Data) und Side-by-Side-Simulationsberichten, mit denen sich ein simuliertes Soll-Szenario unmittelbar gegen den tatsächlichen Prozessverlauf stellen lässt \parencite[Abschnitt~3.1.16 und 3.4, S.~18 und 27--33]{sap2022samplecontent}. Auch die Fallstudie von Adolfsson/Saleh benennt die Verbindung rückblickender Analyse mit vorausschauendem Process Mining und \gls{rpa} als zentrales künftiges Feld \parencite[S.~73 f.]{adolfsson2026o2c}, und \gls{ki}-gestützte Auftragssteuerung gilt in der \gls{o2c}-Literatur zu \gls{sap} \gls{sd} als nächste Ausbaustufe \parencite[S.~111 f.]{nadukuru2024o2c}. Schließlich senken quelloffene Werkzeuge wie \gls{pm4py} die Einstiegshürde, datengetriebene Prozessanalyse ist damit nicht mehr Großkonzernen vorbehalten. Ein klar umrissener, wertschöpfungsnaher Prozess wie Order-to-Cash ist dafür der richtige erste Anwendungsfall.

Insgesamt zeigt die Arbeit einen methodischen Mehrwert gegenüber der klassischen, workshopbasierten Fit/Gap-Analyse. Deren Ergebnis hängt von der Vollständigkeit der Dokumentation und der Erinnerung der Beteiligten ab, während Process Mining sich auf die objektive Aufzeichnung des \gls{erp}-Systems stützt und die Lücken zählt statt schätzt. Für ein Unternehmen bedeutet das nicht nur eine genauere Diagnose, sondern auch eine dauerhaft wiederholbare: Dieselbe Auswertung kann nach jeder Prozessänderung erneut ausgeführt werden und macht so den Erfolg der hier abgeleiteten Maßnahmen über die Zeit messbar. Damit schließt sich der Bogen von der Datenstruktur über den Algorithmus bis zur betrieblichen Steuerung, den diese Arbeit am Beispiel des Order-to-Cash-Prozesses nachgezeichnet hat.

Für die betriebliche Praxis bedeutet dies, dass der Einstieg in Process Mining keine große Vorabinvestition erfordert. Notwendig sind ein klar abgegrenzter Prozess, ein Verständnis der zugrunde liegenden Belegtabellen und die Bereitschaft, die eigenen Abläufe an gemessenen statt an vermuteten Zahlen zu bewerten. Der \gls{o2c}-Prozess bietet dafür ideale Voraussetzungen, weil er hohe Fallzahlen, eine klare Belegstruktur und eine unmittelbare finanzielle Wirkung verbindet. Gelingt der Nachweis eines Nutzens an diesem Prozess, ist der Weg zu weiteren Anwendungsfeldern methodisch bereits geebnet.
